% TOP_PROBLEM: {"id": 1500021, "title": "Algebraic Stories — Polynomial generation", "category_id": 4, "category": {"id": 4, "name": "algebra", "display_name": "Algebra", "description": "Group theory, ring theory, field theory, and algebraic structures.", "slug": "algebra", "order_index": 4, "created_at": "2026-07-31T15:26:25.671Z"}, "status": "partially_solved"}
% FIRST_POSTED: null
% ENTRY_KIND: statement_only
% Imported from: batch06.tex; UnsolvedMath ID 1500021
% Compile twice: lualatex 1500021.tex
\documentclass[11pt]{article}
\usepackage{fontspec}
\setmainfont{Latin Modern Roman}
\usepackage{amsmath,amssymb,mathtools,unicode-math}
\setmathfont{Latin Modern Math}
\usepackage{xurl,hyperref}
\hypersetup{colorlinks=true,urlcolor=blue,pdftitle={UnsolvedMath Problem 1500021}}
\newcommand{\sref}[2]{\hyperlink{source:#1:#2}{[S#2]}}
\newcommand{\cataloguescope}[1]{\paragraph{Mathematical question.}#1}
\begin{document}
% UnsolvedMath ID 1500021; source catalogue code AMR-014-0021
\setcounter{section}{1500020}
\section{Algebraic Stories — Polynomial generation}\label{1500021}
{\small\sffamily UnsolvedMath ID 1500021 · Algebra}
\subsection{Definitions and mathematical statement}

\paragraph{Shared notation.}$\mathbb N=\{1,2,\ldots\}$, and $\mathbb Z,\mathbb Q,\mathbb R,\mathbb C$ denote the integers, rationals, reals and complex numbers. Any other conventions are specified locally.


Let $p$ be a prime and identify each element $a\in\mathbb F_p$ with its integer representative $\bar a\in\{0,\ldots,p-1\}$ when used as an exponent. Define
\[Z(f)=\{a\in\mathbb F_p:f(a)=0\},\qquad\phi(f)=\sum_{a\in Z(f)}x^{\bar a},\quad f\in\mathbb F_p[x].
\]
After one application the degree is at most $p-1$, so every orbit eventually enters a finite periodic cycle. A cycle of length $q\geq1$ consists of distinct polynomials $f_0,\ldots,f_{q-1}$ with $\phi(f_j)=f_{j+1}$, indices taken modulo $q$.
\paragraph{Statement recorded in the supplied source.}
For each prime $p$, determine all periodic cycles of $\phi$ in one variable and, in particular, all their possible lengths. \sref{1500021}{2}
\subsection{Short English statement}
Classify the cycles and cycle lengths of the map that records a polynomial’s zeros as monomial exponents.
\subsection{Sources}
\begin{description}
\item[\hypertarget{source:1500021:1}{S1}] ULAM.AI and the UnsolvedMath Contributors, UnsolvedMath: A Curated Collection of Open Mathematics Problems, record 1500021 (AMR-014-0021). CC BY 4.0. \url{https://huggingface.co/datasets/ulamai/UnsolvedMath}.
\item[\hypertarget{source:1500021:2}{S2}] R. Fröberg, S. Lundqvist, A. Oneto and B. Shapiro, \emph{Algebraic Stories from One and from the Other Pockets}, arXiv:1801.01692 (2018), §4.3, Problem N. \url{https://arxiv.org/abs/1801.01692}.
\end{description}
\label{end:1500021}
\end{document}
